%% kapa: exam=1 round=37 year=2026 period=2 form=A subject=accounting subject_ko=회계학 seq=10 number=50
%% source: Q-Net 감정평가사 기출문제 (cst00309) artlSeq=5258924 file=raw/1008/5258924/2026년 제37회 감정평가사 1차 2교시.pdf
%% conversion: kordoc 4.24.51 -> tools/gen1.py (auto); review=proofread-sonnet; flags=none
\begin{question}{50}
(주)감평은 20×1년 초 (주)대한이 발행한 주식을 취득하고 이를 당기손익-공정가치 측정 금융자산으로 분류하였으며, 관련 자료는 다음과 같다. 동 금융자산에 대한 회계처리가 (주)감평의 20×1년도 당기순이익에 미치는 영향은? (단, 금융자산의 단가 계산 시 이동평균법을 적용한다.)
\begin{framedtext}
○ 20×1년 1월 1일 (주)대한이 발행한 주식 400주(주당 액면금액 ￦5,000)를 주당 ￦10,000(공정가치)에 현금으로 취득하고, 주당 ￦100의 수수료를 별도로 지급함 \par
○ 20×1년 9월 1일 (주)대한은 무상증자를 결의하였으며, (주)감평은 이로 인하여 동 일자에 (주)대한 주식 100주를 수령하였음(당일 주당 공정가치 ￦11,000) \par
○ 20×1년 12월 1일 보유하고 있던 (주)대한의 주식 60 \%를 주당 현금 ￦12,000에 처분하였으며, 수수료로 주당 ￦120을 별도로 지급함 \par
○ 20×1년 12월 31일 (주)대한의 주식 1주의 공정가치가 ￦13,000이고, 예상 처분 수수료는 주당 ￦130임 \par
\end{framedtext}
\begin{choices}
  \item ￦2,014,000 증가
  \item ￦2,100,000 증가
  \item ￦2,124,000 증가
  \item ￦2,200,000 증가
  \item ￦2,224,000 증가
\end{choices}
\end{question}
\answer{3}
